All experiments use $N=50$ bit positions with a voting portion of $0.5$
($N_s = N_t = 25$), a population of $M=300$ individuals, $V=2$ bits voted on per
round, and 100 voting iterations per run, averaged over 50 independent runs per
configuration. All individuals share identical initial voting bits but have
unique, randomly generated individual bits.

\subsection{Experiment 1: $K \times \alpha$ Sweep (Direct Democracy)}

We sweep landscape ruggedness $K \in \{0, 1, 5, 10, 15, 20\}$ and
cross-dependency $\alpha \in \{0.0, 0.25, 0.5, 0.75, 1.0\}$ under direct
democracy for all eight voting methods. For each configuration we record the mean
fitness, fitness variance, minimum and maximum fitness, and the full fitness
distribution at every iteration.

This sweep addresses the core question: which voting methods are most effective
at collective optimization, and how does this depend on the complexity of the
landscape ($K$) and the degree to which shared decisions interact with individual
heterogeneity ($\alpha$)?

\subsection{Experiment 2: Representative Democracy}

Using a subset of the landscape configurations ($K \in \{5, 10, 15\}$, same
$\alpha$ grid), we compare direct democracy to representative democracy with
$C \in \{3, 5\}$ candidates and selection temperatures
$\tau \in \{\infty\text{ (uniform)}, 1.0, 0.1\}$. All eight voting methods are
tested in both the direct and representative settings.

This experiment examines how the introduction of representation---and the
mechanism by which candidates are selected---affects optimization quality and
distributional equity.

% TODO: Adjust candidate counts and temperatures after initial results if needed.
